\section{Introduction}
\label{sec:intro}
Mutation testing is a technique used to evaluate a test suite by creating mutants and testing whether the test suite can recognize them. However, mutation testing often creates an enormous number of mutants. This huge number of mutants significantly increases the practical cost and makes mutation testing hard to apply in practice. This financial cost is one reason for a long line of research related to mutant reduction/selection that has existed for years \cite{guizzo2022sentinel, zhang2019empirical, ojdanic2023evolving}. Choosing a small mutant subset that is "worth running" would provide practical value for developers, CI pipelines, or even at industrial scale.

Recent studies have brought LLMs into mutation testing in two main ways. The first way is generating semantic mutants: Wang et al.'s comprehensive research on 851 real bugs in Java (Defects4J 2.0 + ConDefects) shows that mutants generated from LLMs reach 76.47\% real-bug detection compared to 44.15\% for rule-based techniques - a 1.75x improvement - even though this comes at the price of a non-compilable rate, duplicate rate, and higher equivalent rate of 32.73 / 9.38 / 3.64 percentage points, respectively \cite{wang2025comprehensive}; after that, SMART raised validity from 42.89\% to 65.6\% and real-bug detection to 92.61\% due to RAG + focused chunking + fine-tuning \cite{wang2026smart}; LLMorpheus generates mutants through prompt templates \cite{tip2025llmorpheus}. The second way is detecting/filtering equivalent mutants with 98\% precision \cite{tanhaei2025gemllm}; Tian et al. evaluate the LLM system's ability to detect equivalent mutants \cite{tian2024equivalent}. Besides that, predictive mutation testing (PMT) trains models to predict the kill/survive result to reduce practical costs: Seshat predicts the kill matrix via a natural language channel in code \cite{kim2022seshat}, MutationBERT improves prediction using CodeBERT \cite{jain2023mutationbert}, and SODA surpasses these baselines with kill-F1 improvements ranging from 5.32\% to 114.92\% depending on configuration \cite{zhao2024soda}. Finally, selective mutation research learns how to build or select valuable mutants: LEAM/LEAM++ learn how to create filtering mutants \cite{tian2022leam, tian2025leamplusplus}, Cerebro chooses subsuming static mutants \cite{garg2023cerebro}, and MuDelta chooses mutants related to a commit \cite{ma2021mudelta}. Overall, machine learning has proven its potential to reduce mutation testing costs across multiple stages.

However, no study has evaluated GPT-4o-mini as a mutant selection decision maker or directly compared it against random fixed-seed selection under an identical mutant budget. The common thread in the above LLM studies is that they use LLMs to create mutants, filter equivalent mutants, and generate tests. There is no research that uses an LLM as a mutant selector to choose a subset and compares it with random selection under the same mutant budget and the same metrics. The closest research that matches this design is Jimenez et al. \cite{jimenez2018naturalness}, who compared naturalness-based selection (n-gram, not LLM) with random selection on 5 Defects4J projects and concluded that naturalness-based selection is only equivalent to (slightly inferior to) random selection, with a small effect size - this leaves the question of whether a modern LLM can perform better than an n-gram language model in the same role of selecting mutants. PRIMG \cite{bouafif2025primg} uses LLM prioritization but on Solidity smart contracts and for test generation, and it still does not address the choose-vs-random problem in Java.

This paper fills this gap through a comparative experiment, pre-registered with a full hypothesis, seed, and criteria before running. We use GPT-4o-mini to rate the usefulness, from 1 to 5, of all 11,040 PIT mutants from three Defects4J projects belonging to three different domains - Commons Codec (encoding), Commons CSV (CSV parsing), and JacksonCore (JSON parsing) - then choose the top 10\% from each (project, version, class) group and compare it with random selection under the same budget and on the same pool. Unlike other studies that generate mutants and have to process non-compilable mutants, our design keeps PIT's valid mutant pool intact, so that all differences in mutation scores can be attributed *entirely* to the quality of the selection decision. Besides the main numbers, we also analyze the two ways of combining valid results (by group and by mutant), leading to two different statistical conclusions on the same data - a phenomenon that, to our knowledge, has not been reported in any article on mutant selection.

To summarize, this paper contributes:
\begin{enumerate}
\item the first empirical evaluation of GPT-4o-mini as a mutant-selection decision-maker against same-budget random selection, using mutation score as the primary metric on 11,040 PIT mutants from three Defects4J projects spanning three domains (encoding, CSV parsing, JSON parsing).
\item Experimental results showing a +4.6-point overall mutation-score advantage that is not uniform across groups (group-weighted p=0.059 versus mutant-weighted Monte-Carlo p=0.0002) - a unit-of-analysis finding that matters for future selection studies.
\item A publicly available replication package - dataset, prompts, pipeline, and raw results.
\end{enumerate}

The rest of this paper is structured as follows: Section 2 reviews related work according to three themes. Section 3 describes the methodology, including the Defects4J dataset and mutant generation using PIT, the mutant selection pipeline for GPT-4o-mini and random fixed-seed selection under the same mutant budget, the mutation score metric, and the statistical tests and thresholds that were pre-registered. Section 4 presents the experimental results and reports the results for each research question. Section 5 discusses the reasons for and meaning of the results, especially the differences between the projects and the method's applicability in practice. Section 6 analyzes threats to validity and the limitations of the experimental design. Section 7 concludes and suggests directions for further research.

